\documentclass{article}
\usepackage[paperwidth=595.3bp,paperheight=841.9bp,margin=0bp]{geometry}
\usepackage{fontspec,xeCJK,amsmath,graphicx,booktabs,eso-pic,fix-cm}
\setmainfont{texgyretermes-regular.otf}[BoldFont=texgyretermes-bold.otf,ItalicFont=texgyretermes-italic.otf]
\setCJKmainfont{FandolSong-Regular.otf}[BoldFont=FandolHei-Regular.otf]
\pagestyle{empty}
\setlength{\parindent}{0pt}
\setlength{\parskip}{3bp}
\setlength{\abovedisplayskip}{3bp}
\setlength{\belowdisplayskip}{3bp}
\setlength{\abovedisplayshortskip}{2bp}
\setlength{\belowdisplayshortskip}{2bp}
\newcommand{\PlaceBlock}[3]{%
  \AddToShipoutPictureFG*{\AtPageLowerLeft{%
    \setlength{\unitlength}{1bp}%
    \put(90,#1){\makebox(0,0)[lt]{\begin{minipage}[t]{415.3bp}%
      \vspace{0pt}\fontsize{#2bp}{13bp}\selectfont #3%
    \end{minipage}}}%
  }}%
}
\input{numbers.tex}
\begin{document}
\PlaceBlock{676}{9.3}{
灯丝电压2.00 V，第一栅压2.10 V；原记录电压符号$V$下文记为$U$。
\par
\includegraphics[width=\linewidth]{curves.pdf}
\vspace{-5bp}
\par
左图为实测点及保形插值连线，峰顶作局部二次拟合；右图为峰位及线性拟合。
每峰取实测最大值点及左右各两点，平顶取先出现的最大值点为中心$U_c$：
\[
 I_A=a(U-U_c)^2+b(U-U_c)+c,\qquad U_n=U_c-\frac{b}{2a}.
\]
\vspace{-4bp}
{\centering
\setlength{\tabcolsep}{6.1bp}
\begin{tabular}{@{}c*{7}{r}@{}}
\toprule
$U_{G2A}/\mathrm{V}$ & $U_1$ & $U_2$ & $U_3$ & $U_4$ & $U_5$ & $U_6$ & $U_7$ \\
\midrule
\PeakRows
\bottomrule
\end{tabular}
\par}
\vspace{7bp}
各峰电压单位为V。逐差法将前3峰与后3峰配对；最小二乘法拟合$U_n=A+nU_0$：
\[
 U_0^{(\text{逐差})}=\frac{(U_5-U_1)+(U_6-U_2)+(U_7-U_3)}{12},\qquad
 U_0^{(\text{拟合})}=\frac{\sum_{n=1}^{7}(n-4)(U_n-\overline U)}
 {\sum_{n=1}^{7}(n-4)^2}.
\]
{\fontsize{8.8bp}{12bp}\selectfont\FitLines\par}
相对误差按指导书公认值11.7 V计算：$\delta=|U_0-11.7|/11.7\times100\%$。
{\centering
\setlength{\tabcolsep}{10bp}
\begin{tabular}{@{}ccccc@{}}
\toprule
$U_{G2A}/\mathrm{V}$ & 逐差法/V & $\delta$/\% & 最小二乘法/V & $\delta$/\% \\
\midrule
\ResultRows
\bottomrule
\end{tabular}
\par}
}
\PlaceBlock{219}{9.5}{
两种方法均得到氩的第一激发电位约11.7 V，与指导书公认值相符。电流随加速电压
反复出现峰谷，表明电子与氩原子发生具有特定能量损失的非弹性碰撞，支持原子能级分立。
拒斥电压由7.20 V降至6.70 V后，同一加速电压下的电流总体增大，峰位略向低电压移动，
而拟合斜率基本不变。
\par
峰间距随序号总体增大，线性斜率反映本测量区间的平均间距。
0.5 V步进及峰顶较平坦限制峰位定位；电子能量分散、碰撞概率随能量变化和电场非均匀性
也会使峰距偏离理想值。局部拟合可细化峰位，但不能消除这些影响。
}
\null\newpage
\PlaceBlock{661}{10.5}{
\textbf{1.}\quad
在本实验条件下，电子达到第一激发能附近便容易发生非弹性碰撞并损失能量，
随后重新加速并重复这一过程，因此曲线主要呈现第一激发能对应的周期。
更高激发通道可能存在，但在本装置条件下难以从曲线中分辨。
\par\vspace{15bp}
\textbf{2.}\quad
同类相邻峰或谷的电压差都近似为第一激发电位；接触电势等引起的固定偏移，
在作差或带截距拟合中可以抵消。峰、谷的绝对位置和形状受拒斥电压及电子能量分布影响，
所以两种取值可能有差别。应选取清晰的同类极值并统一拟合方法；平顶的峰位定位较敏感，
不能将相邻的一个峰与一个谷相减作为第一激发电位。
\par\vspace{15bp}
\textbf{3.}\quad
加速电场增强时，空间电荷对电子输运的限制减弱，板极收集电流的背景升高。
电子初速度以及碰撞位置、次数存在分散，处于波谷时仍有一部分电子能克服拒斥电场到达板极，
因此谷值电流不为零，并随加速电压增大而总体抬升。
}
\null\newpage
\end{document}
